\section{Method}
\label{sec:method}

\subsection{The bandlimited-Laplacian SPDE and its mode-space solution}
\label{sec:method:spde}

We consider the linear bandlimited-Laplacian SPDE \eqref{eq:lin-spde}
on $\mathbb{T}^d \times [0, T]$, where the noise is bandlimited:
$\dot W_{K_{\max}}(x, t) = \sum_{k \in \mathcal{H}_{K_{\max}}}
[\dot B_k^a(t)\cos(k \cdot x) + \dot B_k^b(t)\sin(k \cdot x)]$,
with $\dot B_k^{a/b}$ independent standard white noises and
$\mathcal{H}_{K_{\max}} = \{k \in \mathbb{Z}^d : 0 < |k| \le K_{\max}\}$
the lex-ordered half-space of integer modes (one representative per
$\cos$/$\sin$ conjugate pair, $k = 0$ excluded since the noise has zero
mean).

By the spectral representation of the heat semigroup, the solution
admits the mode-space decomposition
$u(x, t) = \sum_{k \in \mathcal{H}} [a_k(t)\cos(k \cdot x)
+ b_k(t)\sin(k \cdot x)]$ where each
$a_k, b_k$ satisfies the per-mode SDE
\begin{equation}
da_k(t) = -|k|^2 a_k(t)\, dt + \sigma\, dB_k^a(t), \qquad
a_k(0) = a_{0,k},
\label{eq:ou-mode}
\end{equation}
i.e., an Ornstein--Uhlenbeck process with rate $|k|^2$, stationary
variance $\sigma^2/(2|k|^2)$, and analogous equation for $b_k$.

This factored structure is the foundation of the mode-space
architecture: \emph{the spatial Laplacian becomes diagonal in the
Fourier basis}, so only the time evolution of each coefficient
$a_k(t)$, $b_k(t)$ remains to be approximated. Autograd Laplacians and
FFTs disappear, leaving an $L^2(0, T)$ problem per mode.

\subsection{Trial space and half-period basis}
\label{sec:method:architecture}

For each mode $k$ we parameterise the trial coefficient as
\begin{equation}
a_{\theta, k}(t) = a_{0, k}\,e^{-|k|^2 t} + \sum_{l = 1}^L
\alpha_{k, l}\,\phi_l(t),
\label{eq:trial-ansatz}
\end{equation}
where $\phi_l(t) = \sqrt{2/T}\sin((2l-1)\pi t/(2T))$ is the half-period
sine basis \eqref{eq:half-period-basis}. The leading term is the
deterministic heat semigroup applied to the initial condition: it
satisfies $\partial_t a + |k|^2 a = 0$ exactly and equals $a_{0, k}$ at
$t = 0$. Because $\phi_l(0) = 0$ for every $l$, the trial automatically
satisfies the initial condition $a_{\theta, k}(0) = a_{0, k}$. The
analogous ansatz parameterises $b_{\theta, k}$.

The half-period basis is the eigenbasis of $-\partial_t^2$ with
Dirichlet at $t = 0$ and Neumann at $t = T$:
$-\phi_l''(t) = \mu_l\,\phi_l(t)$ with eigenvalues
$\mu_l = ((2l-1)\pi/(2T))^2$. The Dirichlet condition at $t = 0$
matches the IC; the Neumann condition at $t = T$ leaves the terminal
value of $\phi_l$ free, with $\phi_l(T) = (-1)^{l-1}$. This is the
crucial property: the trial space can represent the OU noise
accumulation at terminal time, $a_{\theta, k}(T) = a_{0, k}e^{-|k|^2 T}
+ \sum_l (-1)^{l-1}\alpha_{k, l}$, with $L$ degrees of freedom in the
terminal value.

\subsection{Galerkin formulation and the closed-form solve}
\label{sec:method:galerkin}

We use the matched Petrov--Galerkin formulation with test space
$\{\phi_{l'}(t)\}_{l' = 1}^L$. For each mode $k$, the residual
projected onto the test space is
\begin{equation}
R_{k, l'}(\theta) = \int_0^T \!\bigl[\partial_t a_{\theta, k}(t)
+ |k|^2 a_{\theta, k}(t)\bigr]\,\phi_{l'}(t)\, dt
- \sigma\,I_{k, l'}^a,
\label{eq:residual-projected}
\end{equation}
where $I_{k, l'}^a = \int_0^T \phi_{l'}(t)\, dB_k^a(t)$ is the
scalar Wiener integral of $\phi_{l'}$ against the OU forcing
realisation. The IC term contributes nothing to the residual: by direct
verification, $\int_0^T [\partial_t(a_{0,k}e^{-|k|^2 t}) + |k|^2
a_{0,k}e^{-|k|^2 t}]\,\phi_{l'}(t)\, dt = 0$ as a strong-form integral
(the IC is in the kernel of $\partial_t + |k|^2\,\mathrm{I}$).

Substituting the trial ansatz \eqref{eq:trial-ansatz} into
\eqref{eq:residual-projected} and using the orthonormality of $\phi_l$,
\begin{equation}
R_{k, l'}(\theta) = \sum_{l = 1}^L M_{l, l'}\,\alpha_{k, l}
+ |k|^2\,\alpha_{k, l'} - \sigma\,I_{k, l'}^a,
\label{eq:residual-galerkin}
\end{equation}
where the Galerkin matrix is
\begin{equation}
M_{l, l'} := \langle \partial_t \phi_l, \phi_{l'} \rangle_{L^2(0, T)}.
\label{eq:M-matrix}
\end{equation}

The matrix $M$ is the half-period analogue of the spectral
$\partial_t$ operator. Its diagonal entries are $M_{l, l} = 1/T$
(the boundary term: $\int_0^T \partial_t\phi_l\,\phi_l\,dt
= \phi_l(T)^2/2 = 1/T$) and off-diagonal entries are explicit
(Lemma~\ref{lem:M-explicit}). Crucially, $M$ has both symmetric and
skew-adjoint components: the symmetric part comes from the
non-vanishing boundary at $t = T$ and is what gives the half-period
architecture its convergence rate.

The Galerkin equation $R_{k, \cdot}(\theta) = 0$ in matrix form is
\begin{equation}
(M^T + |k|^2 \mathrm{I})\,\alpha_k = \sigma\,I_k^a,
\label{eq:galerkin-system}
\end{equation}
a linear system in $L$ unknowns per mode $k$, with the same equation
for $\beta_k$ replacing $I_k^a$ by $I_k^b$. The solution
$\alpha_k^* = (M^T + |k|^2\mathrm{I})^{-1}\sigma I_k^a$ is the
closed-form Galerkin coefficient for the linear SPDE. For the linear
case, no PINN training is needed: the spectral coefficients are
obtainable in time $O(JL^3)$.

\subsection{The closed-form Galerkin matrix $M$}
\label{sec:method:M}

\begin{lemma}[Half-period $\partial_t$ matrix]
\label{lem:M-explicit}
The matrix $M_{l, l'} = \langle \partial_t \phi_l, \phi_{l'}
\rangle_{L^2(0, T)}$ for $\phi_l(t) = \sqrt{2/T}\sin((2l-1)\pi t/(2T))$
on $[0, T]$ has the following structure:
\begin{enumerate}
\item Diagonal: $M_{l, l} = 1/T$ for all $l \ge 1$.
\item Symmetric component: $M_{l, l'} + M_{l', l} = (2/T)(-1)^{l + l'}$
for all $l \ne l'$.
\item Closed form (off-diagonal, $l \ne l'$):
$$
M_{l, l'} = \frac{2l-1}{2T}\left[\frac{1+(-1)^{l+l'}}{l+l'-1}
+ \frac{1-(-1)^{l+l'}}{l'-l}\right].
$$
Equivalently: $M_{l,l'} = (2l-1)/((l+l'-1)T)$ when $l+l'$ is even,
and $M_{l,l'} = (2l-1)/((l'-l)T)$ when $l+l'$ is odd.
\end{enumerate}
\end{lemma}

In particular, $M$ is dense (no sparsity to exploit beyond the
diagonal) but is $L \times L$ with $L$ moderate, so the linear solve
\eqref{eq:galerkin-system} is fast at every $d$.

\subsection{Computational summary}
\label{sec:method:summary}

The complete linear-case algorithm:

\begin{algorithm}[h]
\caption{Closed-form half-period Galerkin solve}
\label{alg:halfperiod-solve}
\begin{algorithmic}[1]
\State Build $\mathcal{H}_{K_{\max}}$ (lex half-space modes, $J = |\mathcal{H}|$).
\State Sample IC and Wiener path: $a_{0,k}, b_{0,k}$ and $\{dB_k^{a/b}\}$.
\State Precompute Wiener integrals $I_{k, l'}^{a/b}$ for $k \in \mathcal{H}$, $l' = 1, \ldots, L$.
\State Precompute $M$ (Lemma~\ref{lem:M-explicit}) and Cholesky/LU of $(M^T + |k|^2 \mathrm{I})$ per mode.
\State For each mode $k$: solve \eqref{eq:galerkin-system} for $\alpha_k^*, \beta_k^*$.
\State Reconstruct $u_\theta(x, t)$ via the trial ansatz.
\end{algorithmic}
\end{algorithm}

Storage: $O(JL)$ for the Wiener integrals, $O(L^2)$ for $M$, $O(JL)$
for the coefficients. Compute: $O(JL^3)$ for the per-mode solves
(dominated by the LU of $(M^T + |k|^2\mathrm{I})$ at each $k$, though
$M$ is shared so amortisation reduces this to $O(L^3 + JL^2)$).

\emph{No autograd. No FFT. No exponential cost in $d$.} The only
$d$-dependent quantity is $J = |\mathcal{H}_{K_{\max}}|$, which grows
polynomially with $d$ at fixed $K_{\max}$; the per-mode cost is
independent of $d$.
